% Plantilla mínima de los boletines IES (pandoc --template).
% Solo carga lo estrictamente necesario por el contenido; la tipografía y la
% geometría las decide LaTeX (tamaño carta, 12pt, márgenes por defecto).

\documentclass[12pt,letterpaper]{article}

% Codificación y acentos (pdflatex)
\usepackage[T1]{fontenc}
\usepackage[utf8]{inputenc}

% Fuente escalable (Latin Modern); evita las bitmap EC de METAFONT
\usepackage{lmodern}

% Español: títulos y pies ("Figura"), guionado
\usepackage[spanish]{babel}

% Secciones sin numeración automática (equivalente a \section*{...}): solo se
% muestran los números que el markdown escribe a mano ("1.", "2.", "A.", ...).
\setcounter{secnumdepth}{-1}

% Los títulos nunca quedan huérfanos al pie de página: reserva espacio antes
% de cada sección y, si no cabe el título + unas líneas, salta a la página
% siguiente. Necesario sobre todo cuando el título precede a una tabla longtable.
\usepackage{needspace}
\usepackage{etoolbox}
\pretocmd{\section}{\needspace{3\baselineskip}}{}{}
\pretocmd{\subsection}{\needspace{3\baselineskip}}{}{}
\pretocmd{\subsubsection}{\needspace{3\baselineskip}}{}{}

% Párrafos: penalizar líneas huérfanas (club) y viudas (widow) al cambiar de
% página: la primera línea de un párrafo no queda sola abajo, ni la última sola arriba.
\clubpenalty=10000
\widowpenalty=10000
\displaywidowpenalty=10000

% Matemáticas (variables como T, I, neq)
\usepackage{amsmath,amssymb}

% Imágenes: los diagramas PDF se ajustan al ancho de caja, sin anchos fijos
\usepackage{graphicx}
\makeatletter
\def\maxwidth{\ifdim\Gin@nat@width>\linewidth\linewidth\else\Gin@nat@width\fi}
\def\maxheight{\ifdim\Gin@nat@height>\textheight\textheight\else\Gin@nat@height\fi}
\makeatother
\setkeys{Gin}{width=\maxwidth,height=\maxheight,keepaspectratio}

% Fondo: membrete a página completa (carta) en TODAS las hojas. La ruta la
% inyecta build-boletin.sh con -V membrete=<ruta absoluta>; si la variable no
% se define, el fondo simplemente no se carga. Ojo: en un template de pandoc
% los signos de dolar se interpretan como variables de template, incluso en los
% comentarios; por eso aqui no debe aparecer ningun signo de dolar suelto
% aparte de la variable membrete del bloque condicional.
\usepackage{background}
% NOTA: hay que dar width Y height explicitos (paperwidth x paperheight). Las
% claves globales de Gin (width=\maxwidth, height=\maxheight, keepaspectratio)
% fijadas mas arriba para los diagramas se aplican por defecto a todo
% \includegraphics; con solo width=\paperwidth el fondo quedaria reducido al
% area de texto (margenes). Con ambas dimensiones explicitas y la misma
% proporcion carta del membrete, llena la pagina sin deformar ni recortar.
\backgroundsetup{
  scale=1,
  angle=0,
  opacity=1,
  contents={\includegraphics[width=\paperwidth,height=\paperheight]{/home/kritt/Git/boletin/.opencode/skills/boletin-latex-pdf/assets/membrete.pdf}}
}

% Tablas (los boletines pueden traer tablas Markdown)
\usepackage{longtable,booktabs,array,calc}

% Pies de figura
\usepackage{caption}
\captionsetup{font=small,labelfont=bf}

% Listas compactas de Markdown
\providecommand{\tightlist}{%
  \setlength{\itemsep}{0pt}\setlength{\parskip}{0pt}}

% Enlaces y destinos (al final, como recomienda hyperref)
\usepackage{hyperref}

\begin{document}

\hypertarget{termodinuxe1mica-financiera}{%
\section{Termodinámica Financiera}\label{termodinuxe1mica-financiera}}

\textbf{Análisis Sistémico de Políticas Económicas bajo la Teoría de
Restricciones (TOC)}

\begin{center}\rule{0.5\linewidth}{0.5pt}\end{center}

\hypertarget{titular-y-resumen-ejecutivo}{%
\subsection{1. Titular y Resumen
Ejecutivo}\label{titular-y-resumen-ejecutivo}}

\textbf{El Verdadero Cuello de Botella del Estado: Entre la Contención
Contable y el Flujo de Throughput Social}

El debate entre el modelo neoliberal y el actual no es una simple
disputa ideológica sobre porcentajes del PIB, sino un choque de
paradigmas sobre cómo se gestiona el sistema económico de una nación.
Mientras el modelo neoliberal priorizó la optimización local del gasto
operativo (\(OE\)) y la contención salarial ---destruyendo el mercado
interno y disparando la deuda pública (\(I\))---, el modelo actual
utiliza el salario mínimo y la inversión social como variables de
sincronización para elevar el Throughput (\(T\)) social y reducir la
pobreza multidimensional.

\begin{center}\rule{0.5\linewidth}{0.5pt}\end{center}

\hypertarget{diagnuxf3stico-del-sistema-seguxfan-toc}{%
\subsection{2. Diagnóstico del Sistema según
TOC}\label{diagnuxf3stico-del-sistema-seguxfan-toc}}

Bajo la Teoría de Restricciones, una economía nacional no es una suma de
balances contables aislados, sino un flujo continuo interconectado. La
restricción histórica de México ha radicado en la subutilización de su
mercado interno debido a la depresión sistemática del poder adquisitivo.

La siguiente Nube de Evaporación expone el conflicto fundamental entre
ambos modelos de gestión gubernamental:

\begin{figure}
\centering
\includegraphics{Diagramas/figura-1.pdf}
\caption{Nube de Evaporación: Conflicto entre la ortodoxia contable y la
expansión del Throughput social}
\end{figure}

\begin{center}\rule{0.5\linewidth}{0.5pt}\end{center}

\hypertarget{dicotomuxeda-clave-1-contabilidad-de-costos-tradicional-vs.-contabilidad-del-throughput-ta}{%
\subsection{3. Dicotomía Clave 1: Contabilidad de Costos Tradicional
vs.~Contabilidad del Throughput
(TA)}\label{dicotomuxeda-clave-1-contabilidad-de-costos-tradicional-vs.-contabilidad-del-throughput-ta}}

\begin{itemize}
\tightlist
\item
  \textbf{Contabilidad de Costos Tradicional (Modelo Neoliberal):}
  Concibe los salarios y el gasto público en salud/educación como puros
  costos variables y de operación (\(OE\)) que deben minimizarse. Bajo
  esta lógica, deprimir los salarios mínimos y recortar el gasto social
  (llevándolo a mínimos históricos como el 2.1\% del PIB en 1996) se
  percibe erróneamente como un ``ahorro'' y una ganancia de eficiencia.
  Sin embargo, ignora que asfixiar el poder adquisitivo colapsa el
  consumo y obliga a rescates financieros masivos (como el Fobaproa) que
  disparan exponencialmente la Inversión/Inventario (\(I\) en forma de
  deuda pública).
\item
  \textbf{Contabilidad del Throughput (Modelo Actual):} Reconoce que el
  salario mínimo no es un costo a minimizar, sino una \textbf{señal de
  sincronización} que acelera la velocidad del flujo monetario y
  transaccional en el mercado interno. El gasto en salud y bienestar no
  drena al Estado, sino que eleva la capacidad productiva y la robustez
  del sistema, maximizando el Throughput (\(T\)) social.
\end{itemize}

\begin{center}\rule{0.5\linewidth}{0.5pt}\end{center}

\hypertarget{dicotomuxeda-clave-2-teoruxeda-monetaria-moderna-mmt-vs.-teoruxeda-de-restricciones-toc}{%
\subsection{4. Dicotomía Clave 2: Teoría Monetaria Moderna (MMT)
vs.~Teoría de Restricciones
(TOC)}\label{dicotomuxeda-clave-2-teoruxeda-monetaria-moderna-mmt-vs.-teoruxeda-de-restricciones-toc}}

\begin{itemize}
\tightlist
\item
  \textbf{Visión MMT:} Analiza la soberanía monetaria enfatizando que un
  Estado emisor de su propia moneda no enfrenta restricciones
  financieras nominales ni de ``bancarrota'', sino límites de recursos
  reales y presión inflacionaria. Sostiene que el gasto público
  deficitario es el motor necesario para proveer liquidez y empleo
  cuando el sector privado desinvierte.
\item
  \textbf{Visión TOC:} Complementa y acota el análisis al recordar que,
  aun con soberanía monetaria, \textbf{el dinero es solo un habilitador
  del flujo físico y productivo}. Si el Estado inyecta recursos o emite
  moneda sin que exista la capacidad estructural real en la restricción
  del sistema (infraestructura, seguridad, cadenas de suministro o
  energía), el resultado no es más valor real, sino inflación y
  acumulación de ruido e inventarios improductivos. El modelo actual
  acierta al ligar la expansión del gasto a la recuperación tangible de
  capacidades de la base (reducción de la pobreza multidimensional),
  evitando que el recurso se disperse en ineficiencias.
\end{itemize}

\begin{center}\rule{0.5\linewidth}{0.5pt}\end{center}

\hypertarget{conclusiuxf3n-y-prescripciuxf3n-estratuxe9gica-5-pasos-de-enfoque-de-goldratt}{%
\subsection{5. Conclusión y Prescripción Estratégica (5 Pasos de Enfoque
de
Goldratt)}\label{conclusiuxf3n-y-prescripciuxf3n-estratuxe9gica-5-pasos-de-enfoque-de-goldratt}}

Aplicando los \textbf{5 Pasos de Enfoque} a la gestión macroeconómica
actual del país:

\begin{enumerate}
\def\labelenumi{\arabic{enumi}.}
\tightlist
\item
  \textbf{Identificar la restricción:} El cuello de botella ya no es
  meramente la falta de poder adquisitivo básico (parcialmente
  aliviado), sino los rezagos estructurales en \textbf{seguridad
  pública, certidumbre jurídica, infraestructura logística y
  productividad sistémica}.
\item
  \textbf{Explotar la restricción:} Maximizar el rendimiento de la
  infraestructura y los programas sociales actuales mediante la
  digitalización, la simplificación regulatoria y la eficiencia en la
  ejecución del gasto público.
\item
  \textbf{Subordinar todo lo demás a la restricción:} Alinear la
  política fiscal, educativa y comercial (incluyendo la revisión del
  T-MEC) para que sirvan de soporte directo a la seguridad territorial y
  al desarrollo de proveeduría nacional.
\item
  \textbf{Elevar la restricción:} Destinar inversión de capital (\(I\))
  estratégica y masiva a romper los cuellos de botella físicos (energía
  limpia, polos de desarrollo regional, educación tecnológica y combate
  frontal a la extorsión y la violencia).
\item
  \textbf{Prevenir la inercia:} Evitar caer en la complacencia de los
  indicadores históricos superados; el entorno global y la
  relocalización de cadenas (\emph{nearshoring}) exigen mover
  permanentemente la atención hacia las nuevas restricciones de
  competitividad internacional.
\end{enumerate}

\end{document}
